% Scheletro di lavoro, non testo definitivo.
% Le coordinate indicate sono piste di lettura da verificare sul testo greco prima della stesura.
% Metodo: testo originale, traduzione, ricostruzione argomentativa, analisi critica.
\chapter{La città, l'educazione e il filosofo}
\label{cap:platone:citta}

% Problema: Perché l'educazione dei custodi è un problema politico?
% Fonti da verificare: Repubblica II–III 369a–417b.
\section{La formazione della città giusta}

\subsection{Città, bisogni e divisione delle funzioni}

\subsection{Musica, ginnastica e formazione dei custodi}

% Problema: In che senso il sapere può legittimare il governo?
% Fonti da verificare: Repubblica V 449a–480a; VI 484a–502c.
\section{Le condizioni del governo filosofico}

\subsection{Famiglia, donne e comunità dei custodi}

\subsection{Il filosofo e il governo della città}

% Problema: Quale itinerario conduce dall'opinione alla conoscenza?
% Fonti da verificare: Repubblica VI 504a–511e; VII 514a–541b; raccordo con cap. 26.
\section{L'educazione filosofica}

\subsection{Il Bene, la linea e la caverna}

\subsection{Matematica, dialettica e formazione dei governanti}

% Problema: Qual è il rapporto tra forme politiche e caratteri dell'anima?
% Fonti da verificare: Repubblica VIII–IX 543a–592b.
\section{I regimi e il problema della vita buona}

\subsection{La successione delle costituzioni}

\subsection{La critica della tirannide e la felicità del giusto}
